% concatenated from main.tex
\documentclass[11pt]{article} 
\usepackage[utf8]{inputenc}
\usepackage[T1]{fontenc}
\usepackage{lmodern}
\usepackage[margin=1in]{geometry}
\usepackage{amsmath,amssymb,mathtools,amsthm}
\usepackage[numbers]{natbib}
\usepackage{hyperref}
\usepackage{enumitem}
\usepackage{parskip}
% Avoid niche fonts/packages when possible for arXiv compatibility.

%%%%%%%%%%%%%%%%%%%%%%%%%%%%%%%%%%%%%%%%%%%%%%%%%%%%%%%%%%%%
%%%%%%%%%%%%%%%%%%%%%% THEOREMS %%%%%%%%%%%%%%%%%%%%%%%%%%%
%%%%%%%%%%%%%%%%%%%%%%%%%%%%%%%%%%%%%%%%%%%%%%%%%%%%%%%%%%%%
\theoremstyle{plain}
\newtheorem{theorem}{Theorem}[section]
\newtheorem{lemma}[theorem]{Lemma}
\newtheorem{proposition}[theorem]{Proposition}
\newtheorem{corollary}[theorem]{Corollary}
\newtheorem{assumption}[theorem]{Assumption}
\theoremstyle{remark}
\newtheorem{remark}[theorem]{Remark}
%%%%%%%%%%%%%%%%%%%%%%%%%%%%%%%%%%%%%%%%%%%%%%%%%%%%%%%%%%%%
%%%%%%%%%%%%%%%%%%%%%% MACROS %%%%%%%%%%%%%%%%%%%%%%%%%%%%%
%%%%%%%%%%%%%%%%%%%%%%%%%%%%%%%%%%%%%%%%%%%%%%%%%%%%%%%%%%%%
\newcommand{\R}{\mathbb{R}}
\newcommand{\E}{\mathbb{E}}
\newcommand{\B}{\mathbb{B}} 
\newcommand{\Prob}{\mathbb{P}}
\newcommand{\norm}[1]{\left\lVert #1 \right\rVert}
\newcommand{\dist}{\operatorname{dist}}
\newcommand{\conv}{\operatorname{conv}}
\newcommand{\Unif}{\operatorname{Unif}}
\newcommand{\ball}{\mathbb{B}}
\newcommand{\sphere}{\mathbb{S}}
\newcommand{\Var}{Var}
\newcommand{\Binom}{\mathrm{Binom}}
\DeclareMathOperator{\sech}{sech}
\DeclareMathOperator{\csch}{csch}
\newcommand{\1}{\mathbf{1}}

% arXiv-friendly: keep this environment lightweight (no tcolorbox).
% It is only used for visually highlighting statements; on arXiv we
% render it as an indented block.

\usepackage[most]{tcolorbox}
\usepackage{authblk}
\tcbset{
  colback=gray!5,
  colframe=gray!60!black,
  boxrule=0.5pt,
  arc=2pt,
  left=4pt,
  right=4pt,
  top=2pt,
  bottom=2pt,
  fonttitle=\bfseries,
  coltitle=black,
}

\newtcolorbox{propbox}[1][]{
  title={#1},
  enhanced,
  sharp corners,
}





%%%%%%%%%%%%%%%%%%%%%%%%%%%%%%%%%%%%%%%%%%%%%%%%%%%%%%%%%%%%
%%%%%%%%%%%%%%%%%%%%%% TITLE %%%%%%%%%%%%%%%%%%%%%%%%%%%%%%
%%%%%%%%%%%%%%%%%%%%%%%%%%%%%%%%%%%%%%%%%%%%%%%%%%%%%%%%%%%%
\title{Nonsmooth Optimization with Zeroth Order Comparison Feedback}
% === Replace the author block below with your final author list/affiliations ===
\author[1]{Taha EL~BAKKALI}
\author[2]{El~Mahdi~Chayti}
\author[1]{Omar~Saadi}

\affil[1]{UM6P College of Computing}
\affil[2]{Machine Learning and Optimization Laboratory (MLO), EPFL}

\date{}


\begin{document}
\maketitle

\begin{abstract}
We study unconstrained optimization problems of nonsmooth, nonconvex Lipschitz functions, using only noisy pairwise comparisons governed by a known link function. Our goal is to compute a $(\delta,\varepsilon)$-Goldstein stationary point. We combine randomized smoothing with a novel unbiased reduction from comparisons to local value differences. By leveraging a Russian-roulette truncation on the Bernoulli-product expansion of the inverse link, we construct an exactly unbiased estimator for directional differences. This estimator has finite expected cost and variance scaling quadratically with the function gap, $\mathcal{O}(B^2)$, under mild conditions. Plugging this into the smoothed gradient identity enables a standard nonconvex SGD analysis, yielding explicit comparison-complexity bounds for common symmetric links such as logistic, probit, and cauchit.
\end{abstract}
\section{Introduction}
Comparison-based optimization addresses black-box settings where algorithms lack access to calibrated function values or gradients, receiving only ordinal feedback: an indication of whether $f(y)$ is preferred to $f(x)$. This model naturally captures dueling bandits, human-in-the-loop tuning, and RLHF-style preference signals (see, e.g., comparison-oracle-based direct alignment for LLMs~\cite{chen2025compo}). In smooth convex and strongly convex regimes, guarantees have been established under noisy comparison oracles~\cite{pmlr-v139-saha21b}, and subsequent work develops more realistic preference models where comparison reliability degrades near ties, while still yielding end-to-end optimization guarantees~\cite{saha2025dueling}. For smooth nonconvex objectives, the deterministic or noiseless setting has been studied using comparison and ranking feedback with stationarity guarantees~\cite{tang2024zerothorder,kadi2025on,bergou2020stochastic}, and recent work also shows that acceleration is possible even when the oracle only compares objective values~\cite{lobanov2024acceleration}; very recently, the smooth stochastic nonconvex case has also been analyzed, providing nonasymptotic convergence to approximate stationarity under noisy comparisons~\cite{bakkali2026noisy} with increasingly realistic noisy oracles.


In stark contrast, the \emph{nonsmooth} setting remains entirely unexplored. We aim to fill this gap by providing the first comparison-based optimization algorithms with convergence guarantees for nonsmooth convex and nonconvex objectives.

Addressing the nonsmooth nonconvex regime requires a departure from classical gradient-based stationarity, which is ill-defined for such problems. A robust alternative is \emph{Goldstein stationarity}~\citep{goldstein1977optimization}. For locally Lipschitz functions, this concept allows one to characterize approximate stationarity using the Clarke subdifferential~\citep{clarke1990optimization} by considering the convex hull of subgradients within a ball of radius~$\delta$ around a point.

Goldstein stationarity has recently emerged as the central solution concept for derivative-free optimization of Lipschitz nonsmooth objectives, particularly under \emph{zeroth-order value-oracle} access. In that setting, state-of-the-art algorithms combine randomized smoothing with two-point differences to form gradient estimators of a smoothed surrogate, offering finite-time guarantees for finding $(\delta,\varepsilon)$-Goldstein stationary points~\citep{lin2022gradient,pmlr-v195-jordan23a}. However, these results crucially rely on the magnitude of function differences. When algorithms are restricted to \emph{pairwise comparisons} that reveal only relative order, the magnitude information is lost, and we are not aware of any existing nonasymptotic guarantees for computing $(\delta,\varepsilon)$-Goldstein stationary points for general Lipschitz nonsmooth functions.

\paragraph{Comparison model and link function.}
We adopt the standard pairwise-comparison oracle model from preference-based optimization. Given two points $x, y \in \mathbb{R}^d$, the oracle returns a binary outcome $Y \in \{0,1\}$ satisfying
$
\Prob\left(Y = 1 \mid x, y\right) = \sigma\!\left(\Delta\right),
$
where $\Delta := f(y) - f(x)$. The event $Y = 1$ indicates that $y$ is preferred over $x$. Throughout, the link $\sigma:\R\to(0,1)$ is assumed known and strictly increasing. Our analysis applies to a broad class of symmetric links (where $\sigma(-t)=1-\sigma(t)$) whose inverse $g=\sigma^{-1}$ admits a Bernoulli-product expansion on the operated probability interval induced by the bounded-gap regime. Concretely, under the local query design $|\Delta|\le B$, we have $p=\sigma(\Delta)\in[p_-,p_+]$ with $p_\pm=\sigma(\pm B)$, and we assume that on $[p_-,p_+]$,
$
g(p)=\sum_{m\ge1} c_m\big(p^m-(1-p)^m\big)
$
for known coefficients $(c_m)_{m\ge1}$, with absolute convergence. This condition is satisfied by several classical choices, including the logistic link (where $c_m=1/m$), as well as probit and cauchit links on any operated interval $[p_-,p_+]$ bounded away from $\{0,1\}$.

\paragraph{Main contributions.}
We provide the first framework for minimizing nonsmooth nonconvex functions using only noisy comparisons. Our specific contributions are as follows:

\begin{itemize}[noitemsep,topsep=0pt]
    \item \textbf{Unbiased reduction to local values.} We develop a general method to recover unbiased estimates of function value differences from binary comparisons. For the class of links described above, we construct an estimator $\widehat\Delta$ by combining (i) Bernoulli-product estimators $A_m-B_m$, where $A_m=\1\{Y_1=\cdots=Y_m=1\}$ and $B_m=\1\{Y_1=\cdots=Y_m=0\}$ are derived from $m$ repeated comparisons, with (ii) a Russian-roulette truncation scheme. This construction preserves unbiasedness while ensuring finite expected cost.

    \item \textbf{Optimal variance scaling.} We prove that under a geometric truncation schedule and mild scaling conditions on the link, our estimator achieves a variance of $\E[\widehat\Delta^{\,2}]=\mathcal{O}(B^2)$. This quadratic scaling with the gap size $B$ is crucial, as it matches the variance profile of value-based oracles and prevents the noise explosion typically associated with inverse link approximations.

    \item \textbf{Convergence to Goldstein stationarity.} By substituting the true function difference $\Delta$ with our unbiased estimator $\widehat\Delta$ in the randomized smoothing scheme, we derive a stochastic gradient estimator for the smoothed function $f_\delta$. Since $\nabla f_\delta(x)\in\partial_\delta f(x)$, this yields a rigorous algorithm with finite expected comparison complexity for reaching $(\delta,\varepsilon)$-Goldstein stationarity, extending standard nonconvex SGD guarantees to the nonsmooth comparison-based setting.
\end{itemize}


\section{$(\delta,\varepsilon)$-Goldstein Stationarity Through Pairwise Comparisons}

\subsection{Goldstein stationarity and random smoothing}
\label{sec:goldstein_smoothing}



\paragraph{Goldstein $\delta$-subdifferential.}
For a locally Lipschitz function $f:\R^d\to\R$, define the Goldstein $\delta$-subdifferential as
$
\partial_\delta f(x)
\;:=\;
\mathrm{conv}\Big(\bigcup_{\|z-x\|\le \delta} \partial^C f(z)\Big),
$
where $\partial^C f(z)$ denotes the Clarke subdifferential and $\mathrm{conv}$ the convex hull.
We say that $x$ is $(\delta,\varepsilon)$-Goldstein stationary if
$\mathrm{dist}\big(0,\partial_\delta f(x)\big)\le \varepsilon$.

\paragraph{Random smoothing.}
Let $V\sim \mathrm{Unif}(\mathbb B)$ be uniform on the unit Euclidean ball
$\mathbb B:=\{v\in\R^d:\|v\|\le 1\}$. For $\delta>0$, define the smoothed function
$
f_\delta(x)
\;:=\;
\E\big[f(x+\delta V)\big].
$

\paragraph{A standard smoothing-based certificate.}
The following lemma is a well-known consequence of random smoothing and basic properties of
the Clarke subdifferential: the gradient of the ball-smoothed function $f_\delta$ belongs to
the Goldstein $\delta$-subdifferential of $f$, and hence directly certifies
$(\delta,\varepsilon)$-Goldstein stationarity whenever $\|\nabla f_\delta(x)\|\le \varepsilon$.
We include it for completeness; see \cite{lin2022gradient}.
\begin{propbox}
\begin{lemma}[Smoothed gradients certify Goldstein stationarity]\label{lem:smoothing_bridge_goldstein}
Assume $f$ is locally Lipschitz. Then $f_\delta$ is differentiable and, for every $x\in\R^d$,
$
\nabla f_\delta(x)\in \partial_\delta f(x).
$
Consequently,
$$
\mathrm{dist}\big(0,\partial_\delta f(x)\big)\le \|\nabla f_\delta(x)\|.
$$
In particular, any point $x$ satisfying $\|\nabla f_\delta(x)\|\le \varepsilon$
is $(\delta,\varepsilon)$-Goldstein stationary.
\end{lemma}
\end{propbox}


\paragraph{How comparisons enter.}
Random smoothing admits a convenient two-point gradient formula: if
$u\sim \mathrm{Unif}(\mathbb S^{d-1})$, then
$
\nabla f_\delta(x)
\;=\;
\frac{d}{2\delta}\,
\E\!\left[\big(f(x+\delta u)-f(x-\delta u)\big)\,u\right].
$
This representation reduces the problem of estimating $\nabla f_\delta(x)$ to estimating the
\emph{directional value difference}
$
\Delta_u \;:=\; f(x+\delta u)-f(x-\delta u).
$
In particular, any estimator $\widehat{\Delta}_u$ satisfying
$\E[\widehat{\Delta}_u\mid u]=\Delta_u$ yields an unbiased gradient estimator
\[
\widehat g(x)
\;:=\;
\frac{d}{2\delta}\,\widehat{\Delta}_u\,u,
\qquad\text{so that}\qquad
\E[\widehat g(x)\mid x]=\nabla f_\delta(x),
\]
and can therefore be plugged into standard stochastic gradient methods for minimizing $f_\delta$. The remaining challenge is to construct $\widehat{\Delta}_u$ using only comparison feedback.

















\subsection{Comparison model and link function}
\label{sec:link_model}

\paragraph{Comparison model.}
We access $f$ only through a pairwise-comparison oracle. Given any two points $x,y\in\R^d$, the
oracle returns a random variable $Y\in\{0,1\}$ such that
\[
\Prob(Y=1 \mid x,y) \;=\; \sigma\!\big(f(y)-f(x)\big),
\]
where $\sigma:\R\to(0,1)$ is a known and strictly increasing \emph{link function}. We interpret
$Y=1$ as ``$y$ is preferred to $x$''. Writing
$
\Delta := f(y)-f(x)$, and $
p := \sigma(\Delta),
$
repeated queries of the same pair yield i.i.d.\ samples
$Y_1,\dots,Y_m \overset{\mathrm{i.i.d.}}{\sim} \mathrm{Bernoulli}(p)$.

\paragraph{Operated probability range.}
Our algorithm queries only local pairs of the form $x-\delta u$ and $x+\delta u$, where
$u\in \mathbb S^{d-1}$. Define the directional gap
$
\Delta_u \;:=\; f(x+\delta u)-f(x-\delta u).
$
By assuming that $f$ is $L$-Lipschitz, then
$
|\Delta_u| \;\le\; 2L\delta.$
Set
$
B \;:=\; 2L\delta.
$
Then for every queried direction $u$, the corresponding success probability
$
p \;=\; \sigma(\Delta_u)
$
is confined to a fixed interval
\[
p \in [p_-,p_+],
\qquad
p_- \;:=\; \sigma(-B),
\qquad
p_+ \;:=\; \sigma(B).
\]
We call the interval $[p_-,p_+]$ the \emph{operated probability range}.






































\section{Logistic link and unbiased gap estimation from comparisons}
\label{sec:logistic_rr}

\paragraph{Comparison oracle (logistic link).}
Fix $x,y\in\R^d$ and denote the gap
$\Delta \;:=\; f(y)-f(x).$
A comparison query at $(x,y)$ returns $Y\in\{0,1\}$ such that
\[
\Prob(Y=1\mid x,y)
\;=\;
\sigma_\tau(\Delta)
\;:=\;
\frac{1}{1+e^{-\Delta/\tau}}
\;=\;
\sigma(\Delta/\tau),
\qquad
\sigma(t):=\frac{1}{1+e^{-t}},
\]
where $\tau>0$ is the temperature (noise scale). We interpret $Y=1$ as ``$y$ is preferred to $x$''.

\paragraph{Bounded gaps via symmetric queries.}
We enforce bounded gaps by querying symmetric pairs. Given a radius $\delta>0$ and $u\in S^{d-1}$, we compare
$
x+\delta u \quad \text{vs.} \quad x-\delta u.
$
If $f$ is $L$-Lipschitz, then
$
\big|f(x+\delta u)-f(x-\delta u)\big|
\;\le\;
2L\delta
\;=:\;
B.
$
Hence, for every queried pair, $|\Delta|\le B$ and the comparison probability
$p:=\sigma_\tau(\Delta)$ satisfies
\[
p\in[p_-,p_+],
\qquad
p_-:=\sigma_\tau(-B)=\sigma(-B/\tau),
\qquad
p_+:=\sigma_\tau(B)=\sigma(B/\tau),
\]
so $p$ is uniformly bounded away from $0$ and $1$ as soon as $B/\tau$ is bounded below.

\paragraph{Temperature scaling.}
Only the normalized gap $\Delta/\tau$ matters. When $\tau\gg B$, we have $|\Delta|/\tau\ll 1$ and thus
$p=\sigma(\Delta/\tau)\approx 1/2$, i.e., comparisons become nearly random.
In what follows it is natural (and mild) to assume $\tau\lesssim B$ (equivalently $B/\tau\gtrsim 1$),
so that comparisons retain a nontrivial signal-to-noise ratio on the range of queried gaps.

\paragraph{Why not use the plug-in logit estimator.}
A natural approach is to estimate $p$ by $\hat p=\frac{1}{m}\sum_{i=1}^m Y_i$ from $m$ repeated comparisons and output
$\tau\log\!\big(\frac{\hat p}{1-\hat p}\big)$. This is problematic because (i) the logit is nonlinear, hence biased,
and (ii) $\hat p\in\{0,1\}$ with nonzero probability, making the logit undefined and producing heavy tails.
We therefore construct an exactly unbiased estimator with controlled moments.

% ------------------------------------------------------------
\subsection{A series representation}

\begin{propbox}
\begin{lemma}[Logit series]\label{lem:series}
For any $p\in(0,1)$,
\[
\log\frac{p}{1-p}
\;=\;
\sum_{m\ge 1}\frac{p^m-(1-p)^m}{m}.
\]
\end{lemma}
\end{propbox}

\begin{proof}
For $|t|<1$, $\log(1-t)=-\sum_{m\ge1} t^m/m$.
Apply with $t=1-p$ and $t=p$, and subtract.
\end{proof}

Combining Lemma~\ref{lem:series} with $\Delta=\tau\log\!\big(\frac{p}{1-p}\big)$ yields
$
\Delta
\;=\;
\tau \sum_{m\ge1}\frac{p^m-(1-p)^m}{m}.
$

% ------------------------------------------------------------
\subsection{Estimating the series terms from comparisons}

Fix $(x,y)$ and repeat the \emph{same} comparison query $m$ times to obtain i.i.d.\ samples
$Y_1,\dots,Y_m\sim\mathrm{Bernoulli}(p)$. Define
$
A_m:=\1\{Y_1=\cdots=Y_m=1\}$
and
$B_m:=\1\{Y_1=\cdots=Y_m=0\}.$
Then $\E[A_m]=p^m$ and $\E[B_m]=(1-p)^m$, so $A_m-B_m$ is an unbiased estimator of $p^m-(1-p)^m$.

% ------------------------------------------------------------
\subsection{Russian-roulette truncation}
\label{sub_trunc}
Let $M\in\{1,2,\dots\}$ be an integer-valued random variable sampled independently of all comparisons, and let
$q_m:=\Prob(M\ge m).$
For each $m\le M$, we use a \emph{fresh block} of $m$ repeated comparisons (independent across different $m$) to form
$A_m$ and $B_m$, and output the Russian-roulette estimator
$
\widehat\Delta
\;:=\;
\tau\sum_{m=1}^{M}\frac{A_m-B_m}{m\,q_m}.
$

\paragraph{Geometric truncation.}
Fix $\beta\in(0,1)$ and sample $M$ with $\Prob(M=m)=(1-\beta)\beta^{m-1}$.
Then $q_m=\Prob(M\ge m)=\beta^{m-1}$.

% ------------------------------------------------------------
\subsection{Unbiasedness, expected cost, and second moment}

\begin{propbox}
\begin{lemma}[Unbiasedness]\label{lem:unbiased}
For any truncation $M$ independent of the comparison outcomes,
\[
\E[\widehat\Delta]=\Delta.
\]
\end{lemma}
\end{propbox}

\begin{proof}
We have
$
\widehat\Delta
=
\tau\sum_{m\ge1}\frac{\1\{M\ge m\}}{m\,q_m}(A_m-B_m).
$
Since $\E[\1\{M\ge m\}/q_m]=1$ and $\E[A_m-B_m]=p^m-(1-p)^m$, we obtain
$
\E[\widehat\Delta]
=
\tau\sum_{m\ge1}\frac{p^m-(1-p)^m}{m}
=
\tau\log\frac{p}{1-p}
=
\Delta.
$
\end{proof}

\begin{propbox}
\begin{lemma}[Expected comparison cost]\label{lem:cost}
Under geometric truncation with parameter $\beta$, the total number of comparisons used to compute $\widehat\Delta$ is
$N=\sum_{m=1}^{M} m = \frac{M(M+1)}{2}$, and
$
\E[N]=\frac{1}{(1-\beta)^2}.
$
\end{lemma}
\end{propbox}

\begin{proof}
For geometric $M$, $\E[M]=1/(1-\beta)$ and $\E[M^2]=(1+\beta)/(1-\beta)^2$. Hence
$\E[M(M+1)/2]=(\E[M^2]+\E[M])/2=1/(1-\beta)^2$.
\end{proof}

\begin{propbox}
\begin{lemma}[Finite second moment]\label{lem:second}
Assume $|\Delta|\le B$ and geometric truncation $q_m=\beta^{m-1}$. Let
\[
\alpha:=\max\{p_+,\,1-p_-\}.
\]
If $\beta>\alpha$, then $\E[\widehat\Delta^{\,2}]<\infty$. Moreover,
\[
\E[\widehat\Delta^{\,2}]
\le
B^2
+
\tau^2 \sum_{m\ge1}\frac{p_+^m+(1-p_-)^m}{m^2\,q_m}
+
\tau^2\left(\sum_{m\ge1}\frac{p_+^m+(1-p_-)^m}{m\,q_m}\right)^2.
\]
\end{lemma}
\end{propbox}

\begin{proof}
Let $Z_m:=(A_m-B_m)/(m q_m)$ and write $\widehat\Delta=\tau\sum_{m\ge1}\1\{M\ge m\}Z_m$.
Using $\E[\widehat\Delta^{\,2}]=\Var(\widehat\Delta)+(\E\widehat\Delta)^2$ and $\Delta^2\le B^2$, it suffices to bound
$\Var(\widehat\Delta)$. Decompose
\[
\Var(\widehat\Delta)
=
\E[\Var(\widehat\Delta\mid M)]
+
\Var(\E[\widehat\Delta\mid M]).
\]
Conditioned on $M$, the blocks across $m$ are independent, hence
\[
\Var(\widehat\Delta\mid M)
=
\tau^2\sum_{m=1}^{M}\Var(Z_m)
\le
\tau^2\sum_{m=1}^{M}\E[Z_m^2].
\]
Taking expectation in $M$ gives
$
\E[\Var(\widehat\Delta\mid M)]
\le
\tau^2\sum_{m\ge1} q_m\,\E[Z_m^2].
$
Since $(A_m-B_m)^2=A_m+B_m$,
\[
\E[Z_m^2]
=
\frac{\E[A_m+B_m]}{m^2 q_m^2}
=
\frac{p^m+(1-p)^m}{m^2 q_m^2}
\le
\frac{p_+^m+(1-p_-)^m}{m^2 q_m^2},
\]
which yields the first series term after multiplying by $q_m$.
For the second term, $\E[\widehat\Delta\mid M]=\tau\sum_{m=1}^{M}\mu_m$ where
$\mu_m=\E[Z_m]=(p^m-(1-p)^m)/(m q_m)$, so
\[
\Var(\E[\widehat\Delta\mid M])
\le
\E[(\E[\widehat\Delta\mid M])^2]
\le
\tau^2\left(\sum_{m\ge1}\frac{p_+^m+(1-p_-)^m}{m q_m}\right)^2.
\]
Finally, $\beta>\alpha$ ensures convergence since $q_m=\beta^{m-1}$.
\end{proof}





\begin{propbox}
\begin{corollary}[Second moment scales as $O(B^2)$]\label{cor:momentB}
Assume $|\Delta|\le B$ and geometric truncation $q_m=\beta^{m-1}$ with $\beta\in(p_+,1)$, where $p_+=\sigma_\tau(B)$.
If moreover $\tau=\mathcal{O}(B)$, then $\E[\widehat\Delta^{\,2}]=\mathcal{O}(B^2)$.
\end{corollary}
\end{propbox}

\begin{proof}
For the logistic link, $p_- = 1-p_+$, hence $p_+^m+(1-p_-)^m=2p_+^m$.
With $q_m=\beta^{m-1}$ and $\beta\in(p_+,1)$, the two series in Lemma~\ref{lem:second} converge to finite constants
depending only on $p_+/\beta<1$. Therefore $\E[\widehat\Delta^{\,2}]\le B^2 + C\tau^2$ for some $C<\infty$.
Since $\tau=\mathcal{O}(B)$, we have $\tau^2=\mathcal{O}(B^2)$, which implies $\E[\widehat\Delta^{\,2}]=\mathcal{O}(B^2)$.
\end{proof}

% ------------------------------------------------------------
\subsection{Plugging $\widehat\Delta$ into two-point smoothing}
\label{sec:plug_smoothing}

\paragraph{Smoothing objective.}
Fix $\delta>0$, we have
\[
\nabla f_\delta(x)
=
\frac{d}{\delta}\,\E_{u\sim \mathrm{Unif}(S^{d-1})}\big[f(x+\delta u)\,u\big]
=
\frac{d}{2\delta}\,\E_{u\sim \mathrm{Unif}(S^{d-1})}\big[(f(x+\delta u)-f(x-\delta u))\,u\big].
\]


\paragraph{Gradient estimator from comparisons.}
Sample $u\sim\mathrm{Unif}(S^{d-1})$ and estimate the symmetric gap
$
\Delta(x,u):=f(x+\delta u)-f(x-\delta u)
$
by $\widehat\Delta(x-\delta u,x+\delta u)$ using the procedure above (with fresh comparisons).
Define
$
G(x;u)
:=
\frac{d}{2\delta}\,\widehat\Delta(x-\delta u,\,x+\delta u)\,u.
$
By Lemma~\ref{lem:unbiased} and the smoothing identity, $\E[G(x;u)\mid x]=\nabla f_\delta(x)$.

Moreover, by Lipschitzness $|\Delta(x,u)|\le 2L\delta=:B$, so Corollary~\ref{cor:momentB} applies whenever
$\tau=\mathcal{O}(B)$, yielding $\E[\widehat\Delta(x-\delta u,x+\delta u)^2]=\mathcal{O}(B^2)=\mathcal{O}(L^2\delta^2)$.
Therefore,
\[
\E\big[\|G(x;u)\|^2 \,\big|\, x\big]
=
\frac{d^2}{4\delta^2}\,\E\big[\widehat\Delta(x-\delta u,x+\delta u)^2 \,\big|\, x\big]
=
\mathcal{O}(d^2 L^2).
\]
Finally, the expected number of comparisons used to form $\widehat\Delta$ is $\E[N]=1/(1-\beta)^2$ by
Lemma~\ref{lem:cost}, giving an explicit per-iteration comparison cost.



\section{General inverse-link expansions compatible with comparisons}
\label{sec:general_link_series}

This section shows that our comparison-only construction is \emph{not tied to the logistic link}.
The only role of the link is to convert the Bernoulli probability $p$ of a comparison outcome into the
latent gap $\Delta=\tau\,g(p)$ via the inverse link $g=\sigma^{-1}$. We explain a general condition on
$g$ under which $\Delta$ admits an \emph{unbiased} estimator using comparisons only.

\paragraph{Scaled comparison model and symmetry.}
Let $\sigma:\R\to(0,1)$ be strictly increasing and define the inverse link
$
g(p):=\sigma^{-1}(p).
$
We work in the scaled model
$
p=\sigma_\tau(\Delta):=\sigma(\Delta/\tau)$, where 
$
\Delta=\tau\,g(p).$ 
We assume the standard symmetry
$
g(1-p)=-g(p)
\,\,\, \Longleftrightarrow \,\,\,
\sigma(-t)=1-\sigma(t),
$
which holds for logistic, probit, and many common symmetric links.

\paragraph{Operated interval (what we control).}
By query design we restrict to gaps $|\Delta|\le B$. Then
$
p=\sigma(\Delta/\tau)\in[p_-,p_+], 
p_-:=\sigma(-B/\tau)$ and  $p_+:=\sigma(B/\tau)$ 
with $0<p_-<p_+<1$. Under symmetry, $p_-=1-p_+$ so the interval is centered at $1/2$.
Define
$
\alpha:=\max\{p_+,\,1-p_-\}\in\Big(\frac12,1\Big).$
Intuitively, $\alpha$ measures how close the operated interval is to the boundary $\{0,1\}$; staying
away from $0$ and $1$ is what allows stable expansions of $g$.

\paragraph{Comparison-friendly basis.}
For $m\ge1$ define the Bernoulli-product polynomials
\[
b_m(p):=p^m-(1-p)^m .
\]
These functions are \emph{exactly} compatible with comparisons: if
$Y_1,\dots,Y_m\stackrel{\mathrm{i.i.d.}}{\sim}\mathrm{Bernoulli}(p)$ then
$
\widehat b_m
:=\1\{Y_1=\cdots=Y_m=1\}-\1\{Y_1=\cdots=Y_m=0\}
\,\,\, \text{satisfies}\,\,\,
\E[\widehat b_m]=b_m(p).
$
Moreover, for $p\in[p_-,p_+]$ we have the uniform bound
$
|b_m(p)|\le p^m+(1-p)^m \le 2\alpha^m.
$ We only need a mild regularity assumption ensuring that $g$ can be expanded on $[p_-,p_+]$ with
geometrically decaying coefficients.

\begin{assumption}[Analytic inverse on the operated interval]\label{ass:general_link_simple}
The inverse link $g=\sigma^{-1}$ is odd around $1/2$ and admits a holomorphic (complex-analytic) extension
to an open neighborhood of the compact interval $[p_-,p_+]$.
\end{assumption}

\begin{propbox}
\begin{proposition}[General-link Bernoulli-product expansion]\label{prop:general_link_simple}
Under Assumption~\ref{ass:general_link_simple}, there exist constants $C>0$ and $\rho\in(0,1)$
(depending only on $g$ and the interval $[p_-,p_+]$) and coefficients $(c_{2k-1})_{k\ge1}$ such that for all
$p\in[p_-,p_+]$,
$$
g(p)=\sum_{k\ge1} c_{2k-1}\,b_{2k-1}(p),
\qquad
|c_{2k-1}|\le C\,\rho^k.
$$
In particular, the above series converges absolutely and uniformly on $[p_-,p_+]$.
\end{proposition}
\end{propbox}

\begin{proof}[Proof sketch]
Since $g$ is holomorphic on a neighborhood of $[p_-,p_+]$, standard approximation theory yields a sequence of
polynomials $(P_r)_{r\ge1}$ such that $\sup_{p\in[p_-,p_+]}|g(p)-P_r(p)|$ decays geometrically in $r$.
By symmetry, $g$ is odd around $1/2$, so we may take $P_r$ to be odd around $1/2$ as well (e.g., by odd symmetrization). For each $r$, the family $\{b_{1},b_{3},\dots,b_{2r-1}\}$ spans the space of polynomials odd around $1/2$
of degree at most $2r-1$, hence there exist coefficients $(c^{(r)}_{2k-1})_{k\le r}$ such that
$P_r(p)=\sum_{k=1}^{r} c^{(r)}_{2k-1}\,b_{2k-1}(p)$.
Taking $r\to\infty$ and using the geometric approximation rate implies that these finite expansions converge
uniformly to $g$ on $[p_-,p_+]$, yielding an infinite series representation in the $b_{2k-1}$ basis.
Finally, analyticity on a neighborhood of $[p_-,p_+]$ implies geometric decay of the resulting coefficients,
i.e., $|c_{2k-1}|\le C\rho^k$ for some $\rho\in(0,1)$.
\end{proof}

\paragraph{Unbiased estimator for $\Delta$.}
Let $Y\sim\mathrm{Bernoulli}(p)$ represent a single comparison outcome.
For each $k\ge1$, draw an independent block of size $2k-1$:
$Y^{(k)}_1,\dots,Y^{(k)}_{2k-1}\stackrel{\mathrm{i.i.d.}}{\sim}\mathrm{Bernoulli}(p)$,
independent across $k$, and define
\[
A_{2k-1}:=\1\{Y^{(k)}_1=\cdots=Y^{(k)}_{2k-1}=1\},
\qquad
B_{2k-1}:=\1\{Y^{(k)}_1=\cdots=Y^{(k)}_{2k-1}=0\}.
\]
Then $\E[A_{2k-1}-B_{2k-1}]=b_{2k-1}(p)$.
Let $M$ be independent of all comparisons with survival probabilities $q_k:=\Prob(M\ge k)$.
Define the Russian-roulette estimator
$$
\widehat\Delta
:=
\tau\sum_{k=1}^{M}\frac{c_{2k-1}}{q_k}\,\big(A_{2k-1}-B_{2k-1}\big).
$$
\begin{propbox}
\begin{lemma}[Unbiasedness]\label{lem:unbiased_simple}
Assume $p\in[p_-,p_+]$ and $\sum_{k\ge1}|c_{2k-1}|\,\alpha^{2k-1}<\infty$.
Then $\E[\widehat\Delta]=\Delta$.
In particular, under Assumption~\ref{ass:general_link_simple} the estimator above is unbiased.
\end{lemma}
\end{propbox}
\begin{proof}
Write the estimator as
$
\widehat\Delta
=
\tau\sum_{k\ge1}\frac{\1\{M\ge k\}}{q_k}\,c_{2k-1}\,\big(A_{2k-1}-B_{2k-1}\big)$, where $ q_k:=\Prob(M\ge k).
$
We first justify exchanging expectation and the infinite sum.

For each $k$, we have $|A_{2k-1}-B_{2k-1}|\le A_{2k-1}+B_{2k-1}$ and
$$
\E[A_{2k-1}+B_{2k-1}]=p^{2k-1}+(1-p)^{2k-1}\le 2\alpha^{2k-1}.
$$
Moreover, $\E\!\left[\frac{\1\{M\ge k\}}{q_k}\right]=1$ since $\E[\1\{M\ge k\}]=q_k$.
Therefore, by Tonelli's theorem,
\begin{align*}
\sum_{k\ge1}\E\!\left[\left|\frac{\1\{M\ge k\}}{q_k}\,c_{2k-1}\,(A_{2k-1}-B_{2k-1})\right|\right]
&\le
\sum_{k\ge1}|c_{2k-1}|\,
\E[A_{2k-1}+B_{2k-1}] \\
&\le
2\sum_{k\ge1}|c_{2k-1}|\,\alpha^{2k-1}
\;<\;\infty.
\end{align*}

Hence we may interchange expectation and summation:
\[
\E[\widehat\Delta]
=
\tau\sum_{k\ge1}c_{2k-1}\,
\E\!\left[\frac{\1\{M\ge k\}}{q_k}\right]\,
\E[A_{2k-1}-B_{2k-1}].
\]
Finally, the blocks are independent and $\E[A_{2k-1}-B_{2k-1}]=b_{2k-1}(p)$, so
\[
\E[\widehat\Delta]
=
\tau\sum_{k\ge1}c_{2k-1}\,b_{2k-1}(p)
=
\tau g(p)
=
\Delta.
\]
\end{proof}


\begin{remark}[Examples: probit and cauchit]\label{rem:general_link_examples_simple}
Assumption~\ref{ass:general_link_simple} holds for many classical symmetric links on any operated interval
$[p_-,p_+]\subset(0,1)$ bounded away from $\{0,1\}$.
\textbf{Probit:} $\sigma(t)=\Phi(t)$, so $g=\Phi^{-1}$ is holomorphic on a neighborhood of any compact
$[p_-,p_+]\subset(0,1)$.
\textbf{Cauchit:} $\sigma(t)=\tfrac12+\tfrac{1}{\pi}\arctan(t)$, so $g(p)=\tan(\pi(p-\tfrac12))$ is meromorphic with poles
at $p\in\mathbb{Z}+\{0,1\}$ and hence holomorphic on a neighborhood of any compact $[p_-,p_+]\subset(0,1)$.
\end{remark}


\begin{propbox}
\begin{lemma}[Second moment control]\label{lem:moment_simple}
Assume $|\Delta|\le B$ so that $p\in[p_-,p_+]$.
Define
$$
S_2:=\sum_{k\ge1}\frac{c_{2k-1}^2}{q_k}\,\alpha^{2k-1},
\qquad
S_1:=\sum_{k\ge1}\frac{|c_{2k-1}|}{\sqrt{q_k}}\,\alpha^{k-\frac12}.
$$
If $S_1<\infty$ and $S_2<\infty$, then
$$
\E[\widehat\Delta^{\,2}]
\;\le\;
B^2 \;+\; 2\tau^2 S_2 \;+\; 2\tau^2 S_1^2.
$$
In particular, if $\tau=\mathcal O(B)$ and $S_1,S_2=\mathcal O(1)$, then $\E[\widehat\Delta^{\,2}]=\mathcal O(B^2)$.
\end{lemma}
\end{propbox}

\begin{proof}
Let
$
X_k:=\frac{c_{2k-1}}{q_k}\,\big(A_{2k-1}-B_{2k-1}\big),
\qquad
\mu_k:=\E[X_k].
$
Then $\widehat\Delta=\tau\sum_{k\ge1}\1\{M\ge k\}X_k$ and $(\E[\widehat\Delta])^2=\Delta^2\le B^2$, hence
$$
\E[\widehat\Delta^{\,2}]\le B^2+\Var(\widehat\Delta).
$$
Conditioned on $M$, the blocks are independent, so
$$
\Var(\widehat\Delta)
=
\E\!\left[\Var(\widehat\Delta\mid M)\right]
+
\Var\!\left(\E[\widehat\Delta\mid M]\right).
$$

For the first term, using $\Var(X_k)\le \E[X_k^2]$ and $\E[\1\{M\ge k\}]=q_k$,
\[
\E\!\left[\Var(\widehat\Delta\mid M)\right]
\le
\tau^2\sum_{k\ge1} q_k\,\E[X_k^2].
\]
Moreover, $(A_{m}-B_{m})^2=A_m+B_m$ and
$\E[A_m+B_m]=p^m+(1-p)^m\le 2\alpha^m$, hence
$$
\E[X_k^2]
=
\frac{c_{2k-1}^2}{q_k^2}\,\E[A_{2k-1}+B_{2k-1}]
\le
\frac{2c_{2k-1}^2}{q_k^2}\,\alpha^{2k-1}.
$$
Therefore
$
\E\!\left[\Var(\widehat\Delta\mid M)\right]
\le
2\tau^2\sum_{k\ge1}\frac{c_{2k-1}^2}{q_k}\,\alpha^{2k-1}
=
2\tau^2 S_2.
$ For the second term, by Cauchy--Schwarz and $\Var(\1\{M\ge k\})\le q_k$,
\[
\Var\!\left(\E[\widehat\Delta\mid M]\right)
=
\tau^2\Var\!\left(\sum_{k\ge1}\1\{M\ge k\}\mu_k\right)
\le
\tau^2\Big(\sum_{k\ge1}\sqrt{q_k}\,|\mu_k|\Big)^2.
\]
Now $\mu_k=\frac{c_{2k-1}}{q_k}\,b_{2k-1}(p)$ and
\[
|b_{2k-1}(p)|
=|\E[A_{2k-1}-B_{2k-1}]|
\le
\E[(A_{2k-1}-B_{2k-1})^2]^{1/2}
=
\E[A_{2k-1}+B_{2k-1}]^{1/2}
\le
\sqrt{2}\,\alpha^{k-\frac12}.
\]
Hence
$
\sqrt{q_k}\,|\mu_k|
\le
\sqrt{2}\,\frac{|c_{2k-1}|}{\sqrt{q_k}}\,\alpha^{k-\frac12},
$
so
$
\Var\!\left(\E[\widehat\Delta\mid M]\right)
\le
2\tau^2\Big(\sum_{k\ge1}\frac{|c_{2k-1}|}{\sqrt{q_k}}\,\alpha^{k-\frac12}\Big)^2
=
2\tau^2 S_1^2.
$
Combining the two bounds yields
$
\Var(\widehat\Delta)\le 2\tau^2 S_2 + 2\tau^2 S_1^2.
$
\end{proof}

\begin{propbox}
\begin{corollary}[Geometric truncation gives a second-moment bound]\label{cor:moment_geometric_simple}
Assume Assumption~\ref{ass:general_link_simple} and $|\Delta|\le B$, so that $p\in[p_-,p_+]$ and $\alpha<1$.
Let $(c_{2k-1})_{k\ge1}$ be the coefficients provided by
Proposition~\ref{prop:general_link_simple}, so that $|c_{2k-1}|\le C\,\rho^k$ for some $C>0$ and $\rho\in(0,1)$.
Let $M$ be geometric with parameter $\beta\in(0,1)$, so that $q_k=\Prob(M\ge k)=\beta^{k-1}$.
If
$$
\beta>\alpha^2\rho^{2},
$$
then $\E[\widehat\Delta^{\,2}]\le B^2 + C_0\tau^2$ for a constant $C_0$ depending only on
$\alpha,\beta,C,\rho$. In particular, if $\tau=\mathcal O(B)$ then
$\E[\widehat\Delta^{\,2}]=\mathcal O(B^2)$.
\end{corollary}
\end{propbox}

\begin{proof}
By Proposition~\ref{prop:general_link_simple} we may take $|c_{2k-1}|\le C\rho^k$.
With $q_k=\beta^{k-1}$, the quantities in Lemma~\ref{lem:moment_simple} satisfy
\[
S_2
=
\sum_{k\ge1}\frac{c_{2k-1}^2}{q_k}\,\alpha^{2k-1}
\le
C^2\,\alpha\sum_{k\ge1}\frac{\rho^{2k}}{\beta^{k-1}}\,\alpha^{2k-2}
=
C^2\,\alpha\,\rho^2\sum_{k\ge1}\Big(\frac{\alpha^2\rho^{2}}{\beta}\Big)^{k-1},
\]
and
$
S_1
=
\sum_{k\ge1}\frac{|c_{2k-1}|}{\sqrt{q_k}}\,\alpha^{k-\frac12}
\le
C\,\alpha^{1/2}\sum_{k\ge1}\frac{\rho^{k}}{\beta^{(k-1)/2}}\,\alpha^{k-1}
=
C\,\alpha^{1/2}\rho\sum_{k\ge1}\Big(\frac{\alpha\rho}{\sqrt{\beta}}\Big)^{k-1}.$
The condition $\beta>\alpha^2\rho^{2}$ implies both geometric series converge, hence $S_1,S_2<\infty$ and in fact
$S_1,S_2=\mathcal O(1)$ with constants depending only on $\alpha,\beta,C,\rho$. Applying Lemma~\ref{lem:moment_simple} yields
$
\E[\widehat\Delta^{\,2}]
\le
B^2 + 2\tau^2 S_2 + 2\tau^2 S_1^2
\le
B^2 + C_0\tau^2
$
for some constant $C_0=C_0(\alpha,\beta,C,\rho)$.
If $\tau=\mathcal O(B)$, then $\E[\widehat\Delta^{\,2}]=\mathcal O(B^2)$.
\end{proof}

\subsection{End-to-end comparison complexity for $(\delta,\varepsilon)$-Goldstein stationarity}
\label{sec:main_theorem}

We now combine the unbiased comparison-based gap estimator from Section~3.5 with a standard
nonconvex SGD guarantee to obtain an explicit \emph{expected comparison-complexity} bound.

\paragraph{Algorithm (Comparison-SGD on $f_\delta$).}
Initialize $x_0\in\R^d$. For $t=0,1,\dots,T-1$:
sample $u_t\sim \mathrm{Unif}(\mathbb{S}^{d-1})$, form an unbiased gap estimate
$\widehat{\Delta}_t := \widehat{\Delta}(x_t-\delta u_t,\;x_t+\delta u_t)$ using the
Russian-roulette procedure (with fresh comparisons), define
$
G_t \;:=\; G(x_t;u_t)\;=\;\frac{d}{2\delta}\,\widehat{\Delta}_t\,u_t,
$
and update $x_{t+1}=x_t-\eta G_t$. Output $x_R$ where $R$ is uniform on $\{0,\dots,T-1\}$
and independent of the algorithmic randomness.


% =========================
% Preamble (add once)
% =========================


\newtcolorbox{propboxwide}{
  enhanced,
  breakable,
  colback=white,
  colframe=black,
  boxrule=0.6pt,
  left=6pt,right=6pt,top=6pt,bottom=6pt,
  width=\linewidth
}


\paragraph{Smoothness of the smoothed objective.}
If $f$ is $L$-Lipschitz and $f_\delta(x)=\E_{v\sim\mathrm{Unif}(\mathbb{B}^d)}[f(x+\delta v)]$,
then $f_\delta$ is differentiable and $L_\delta$-smooth with $L_\delta = L\sqrt d/\delta$.

\begin{propbox}
\begin{lemma}[Nonconvex SGD with unbiased gradients]
\label{lem:sgd_generic}
Let $F:\R^d\to\R$ be $L_F$-smooth and bounded below by $F^\star$.
Consider iterates $x_{t+1}=x_t-\eta G_t$ such that, for all $t$,
$
\E\!\left[G_t\,\middle|\,x_t\right] \;=\; \nabla F(x_t)$ and
$\E\!\left[\|G_t\|^2\,\middle|\,x_t\right] \;\le\; V,$
for some finite constant $V$. If $\eta\le 1/L_F$, then
$
\frac{1}{T}\sum_{t=0}^{T-1}\E\big[\|\nabla F(x_t)\|^2\big]
\;\le\; \frac{2(F(x_0)-F^\star)}{\eta T} \;+\; \eta L_F V.
$

\end{lemma}
\end{propbox}
\par\smallskip

\paragraph{Main theorem (expected comparison complexity).}
We instantiate Lemma~\ref{lem:sgd_generic} with $F=f_\delta$ and the estimator $G_t$ above.
By Section~3.5, $\E[G_t\mid x_t]=\nabla f_\delta(x_t)$.
Assume moreover that the gap estimator has a uniform conditional second-moment bound:
\begin{equation}
\label{eq:gap_second_moment}
\E\!\left[\widehat{\Delta}_t^{\,2}\,\middle|\,x_t,u_t\right] \;\le\; C_\Delta\,B^2
\qquad\text{for all }t,
\end{equation}
where $B:=2L\delta$ (since $|\Delta(x_t,u_t)|\le 2L\delta$ for an $L$-Lipschitz $f$)
and $C_\Delta<\infty$ is a constant.

\leavevmode\par\smallskip
\begin{propbox}
\begin{theorem}[Complexity bound for $(\delta,\varepsilon)$-Goldstein stationarity]
\label{thm:comparison_complexity}
Assume $f$ is bounded below and $L$-Lipschitz, and fix $\delta>0$.
Run Comparison-SGD on $f_\delta$ for $T$ steps with step size
$
\eta \;\le\; \frac{1}{L_\delta}\;=\;\frac{\delta}{L\sqrt d},
$
and suppose \eqref{eq:gap_second_moment} holds for some $C_\Delta<\infty$.
Let $f_\delta^\star:=\inf_x f_\delta(x)$ and $\Delta_0 := f_\delta(x_0)-f_\delta^\star$.
Then for $R\sim\mathrm{Unif}\{0,\dots,T-1\}$ independent,
\[
\E\big[\|\nabla f_\delta(x_R)\|^2\big]
\;\le\;
\frac{2\Delta_0}{\eta T} \;+\; \eta L_\delta \cdot
\frac{d^2}{4\delta^2}\,C_\Delta B^2.
\]
In particular, choosing
$
\eta \;=\; \min\!\left\{\frac{1}{L_\delta},\;
\sqrt{\frac{2\Delta_0}{L_\delta T \cdot \frac{d^2}{4\delta^2}C_\Delta B^2}}\right\}
$
yields
\[
\E\big[\|\nabla f_\delta(x_R)\|^2\big]
\;\le\;
2\sqrt{\frac{2L_\delta \Delta_0}{T}\cdot
\frac{d^2}{4\delta^2}C_\Delta B^2}\,.
\]
Consequently, it suffices to take
$
T \;\ge\;
\frac{8L_\delta \Delta_0}{\varepsilon^4}\cdot
\frac{d^2}{4\delta^2}C_\Delta B^2$
to ensure 
$\E[\|\nabla f_\delta(x_R)\|]\le \varepsilon$.
By Lemma~\ref{lem:smoothing_bridge_goldstein}, this implies that $x_R$ is $(\delta,\varepsilon)$-Goldstein stationary.

Moreover, if each gap estimate uses a (random) number $N$ of comparisons with $\E[N]<\infty$,
then the expected total number of comparisons used by the algorithm is at most
\[
\E[\#\text{comparisons}] \;\le\; T\cdot \E[N].
\]
\end{theorem}
\end{propbox}
\par\smallskip

\paragraph{Explicit specialization under $L_\delta=L\sqrt d/\delta$ and geometric truncation.}
Under \eqref{eq:gap_second_moment} with $B=2L\delta$,
\[
\E\!\left[\|G_t\|^2 \,\middle|\, x_t,u_t\right]
\;=\;
\frac{d^2}{4\delta^2}\E\!\left[\widehat{\Delta}_t^{\,2}\,\middle|\,x_t,u_t\right]
\;\le\;
\frac{d^2}{4\delta^2}\,C_\Delta(2L\delta)^2
\;=\;
C_\Delta d^2 L^2,
\]
so Theorem~\ref{thm:comparison_complexity} gives
$
T \;\ge\; \frac{8L_\delta \Delta_0}{\varepsilon^4}\cdot C_\Delta d^2L^2
\;=\;
\frac{8\,C_\Delta\, d^{5/2} L^3\,\Delta_0}{\delta\,\varepsilon^4}.
$
For the logistic construction with geometric truncation parameter $\beta$,
Section~\ref{sub_trunc} gives $\E[N]=1/(1-\beta)^2$, hence the expected comparison complexity is
$
\E[\#\text{comparisons}]
\;\le\;
T\cdot\E[N]
\;=\;
O\!\left(
\frac{C_\Delta\, d^{5/2} L^3\,\Delta_0}{\delta\,\varepsilon^4}\cdot
\frac{1}{(1-\beta)^2}
\right).$

\section{Related Work}

\textbf{Comparison-Based Optimization.}
Optimization utilizing only relative preferences has been extensively studied under the umbrella of dueling bandits and derivative-free optimization. In the \emph{convex} setting, theoretical guarantees are well-established for both noiseless and noisy comparison oracles~\cite{pmlr-v139-saha21b,saha2025dueling}. For \emph{smooth nonconvex} functions, earlier works focused on deterministic feedback or asymptotic convergence~\cite{bergou2020stochastic}. More recently, \cite{tang2024zerothorder} and \cite{kadi2025on} derived non-asymptotic rates for smooth nonconvex optimization using ranking and comparison feedback, respectively. Most recently, \cite{bakkali2026noisy} extended these results to the stochastic regime with noisy comparisons. However, all these works rely on the smoothness of the underlying objective to control the linearization error of the gradient estimators. Our work departs from this by addressing the \emph{nonsmooth} regime, where gradients may not exist and standard smooth comparison-based estimators fail.

\textbf{Nonsmooth Zeroth-Order Optimization.}
For nonsmooth Lipschitz functions, the standard solution concept is $(\delta, \varepsilon)$-Goldstein stationarity~\citep{goldstein1977optimization}. In the setting of \emph{value-based} zeroth-order optimization (where algorithms observe $f(x)$), \cite{lin2022gradient} pioneered the use of randomized smoothing to convert the nonsmooth problem into a smooth surrogate problem, proving convergence for deterministic oracles. This was extended to stochastic value oracles by \cite{pmlr-v195-jordan23a}. Crucially, \textbf{\cite{kornowski2023optimal} established the optimal dimension-dependence for this setting}, refuting earlier conjectures regarding the necessary scaling with dimension. On the fundamental limits of this regime, \textbf{\cite{zhang2020complexity} established tight lower bounds}, confirming the hardness of the problem and the necessity of the Goldstein stationarity relaxation. Our approach leverages the randomized smoothing geometry used in these works but fundamentally differs in the estimation step: we cannot observe function differences directly. Instead, we introduce a novel unbiased estimator for the inverse link function to bridge the gap between binary comparison feedback and the continuous value differences required by the Goldstein stationary framework.
\section{Conclusion and Discussion}

In this work, we presented the first comparison-based optimization algorithm with non-asymptotic convergence guarantees for nonsmooth, nonconvex Lipschitz functions. By targeting $(\delta, \epsilon)$-Goldstein stationarity, we circumvented the ill-defined nature of gradients in the nonsmooth regime. Our key technical innovation is an unbiased estimator for local function differences that operates solely on binary comparison feedback. Through a Russian-roulette truncation of the inverse link's Bernoulli expansion, we achieved an estimator with finite expected cost and variance scaling as $\mathcal{O}(B^2)$. This variance control allows our method to match the convergence rate of value-based zeroth-order methods (up to constants depending on dimension and link parameters), scaling as $\mathcal{O}(d^{5/2}\delta^{-1}\epsilon^{-4})$.

\textbf{Limitations and Future Work.}
Our current analysis relies on the knowledge of the link function $\sigma$ to construct the inverse expansion. Extending this to \emph{unknown} or semi-parametric link models is an important open direction. Additionally, while the $\epsilon^{-4}$ dependence is standard for SGD on smooth surrogates, acceleration techniques for Goldstein stationarity remains a challenging frontier. Finally, our method assumes unconstrained optimization on $\mathbb{R}^d$; adapting the randomized smoothing and comparison geometry to constrained settings or Riemannian manifolds would broaden the applicability of these results to complex control and ranking problems.
\bibliographystyle{plainnat}
\bibliography{references}


\end{document}
